\section{The classification theorem}

Iterating the lifted successor through one visible circuit gives
\[
T_v^n(i,\epsilon)=(i,\epsilon+\hol(v)).
\]
This identity contains the entire classification.

\begin{theorem}[Binary cycle-lift classification]
For every (n\geq3), deterministic (Ctwo)-covariant two-sheeted lifts of a
directed visible (n)-cycle have exactly two fiber-preserving gauge classes.
\begin{enumerate}
\item If (hol=0), the lifted successor has two orbits of length (n), so
the lift is (C_n\sqcup C_n).
\item If (hol=1), the lifted successor has one orbit of length (2n), so
the lift is (C_{2n}).
\end{enumerate}
Each class is unique up to fiber-preserving gauge and directed-cover
isomorphism.
\end{theorem}

\begin{proof}
The preceding propositions give exactly two gauge classes, indexed by
(hol\in\Ztwo). If (hol=0), then (T_v^n) fixes every lifted state. Since
the visible coordinate first returns after (n) steps, each lifted orbit has
length (n), producing two components.

If (hol=1), then (T_v^n=\deck), so no lifted state returns after (n)
steps. A second circuit gives (T_v^{2n}=\deck^2=1). Hence every lifted orbit
has length (2n). Because there are exactly (2n) lifted states, the lift is
one connected cycle. Completeness of holonomy as a gauge invariant proves
uniqueness in both cases.
\end{proof}

